\chapter{Communication Complexity --- Det Basic}
%%% generated by the blueprint pipeline from TCSlib.CommunicationComplexity.DeterministicCC.DetBasic
\section{Overview}

This module formalizes deterministic two-party communication protocols in which Alice holds
an input $x : X$ and Bob holds an input $y : Y$.  It defines the inductive type
\texttt{CommunicationComplexity.Deterministic.Protocol} together with its execution semantics,
worst-case bit complexity, equivalence and correctness predicates, and structural operations
(role-swap and input pull-back).

%%%%%%%%%%%%%%%%%%%%%%%%%%%%%%%%%%%%%%%%%%%%%%%%%%%%%%%%%%%%%%%%%%%%%%%%%%%%%%%
\section{Declarations}

\begin{definition}[Deterministic communication protocol]
\lean{CommunicationComplexity.Deterministic.Protocol}
\statementsource{roughgarden-cc-bootcamp}{
  pdf-d54d402b16e2-p002-b002,
  pdf-d54d402b16e2-p004-b005
}
\statementsource{raoyehudayoff-ch01-deterministic-protocols}{pdf-fa052d8c95c0-p003-b001}
\label{CommunicationComplexity.Deterministic.Protocol}
\leanok
A deterministic two-party communication protocol over input types $X$ (Alice) and $Y$ (Bob)
producing a value of type $\alpha$ is an inductive type with three constructors: an
\emph{output} node carrying the final value, an \emph{alice} node in which Alice applies a
function $f : X \to \mathrm{Bool}$ to her input and branches on the resulting bit, and a
\emph{bob} node in which Bob does the same with a function $g : Y \to \mathrm{Bool}$.
\end{definition}

\begin{definition}[Protocol execution]
\lean{CommunicationComplexity.Deterministic.Protocol.run}
\statementsource{roughgarden-cc-bootcamp}{
  pdf-d54d402b16e2-p002-b002,
  pdf-d54d402b16e2-p004-b005
}
\statementsource{raoyehudayoff-ch01-deterministic-protocols}{pdf-fa052d8c95c0-p003-b002}
\label{CommunicationComplexity.Deterministic.Protocol.run}
\leanok
\uses{CommunicationComplexity.Deterministic.Protocol}
Given a protocol $p$, \texttt{run} evaluates $p$ on inputs $x : X$ and $y : Y$ by
recursively following the bit-branching structure until an output node is reached,
returning the stored value of type $\alpha$.
\end{definition}

\begin{definition}[Communication complexity]
\lean{CommunicationComplexity.Deterministic.Protocol.complexity}
\statementsource{roughgarden-cc-bootcamp}{pdf-d54d402b16e2-p002-b002}
\statementsource{raoyehudayoff-ch01-deterministic-protocols}{pdf-fa052d8c95c0-p003-b003}
\label{CommunicationComplexity.Deterministic.Protocol.complexity}
\leanok
\uses{CommunicationComplexity.Deterministic.Protocol}
The communication complexity of a protocol $p$ is the worst-case total number of bits
exchanged: an output node costs $0$, while an alice or bob node costs $1$ plus the
maximum complexity of the two sub-protocols reached by the bit $\mathtt{false}$ and the
bit $\mathtt{true}$.
\end{definition}

\begin{definition}[Protocol equivalence]
\lean{CommunicationComplexity.Deterministic.Protocol.Equiv}
\label{CommunicationComplexity.Deterministic.Protocol.Equiv}
\leanok
\uses{CommunicationComplexity.Deterministic.Protocol, CommunicationComplexity.Deterministic.Protocol.run}
Two protocols $p$ and $q$ over the same input and output types are \emph{equivalent} if
$\texttt{run}\,p = \texttt{run}\,q$, i.e.\ they produce identical outputs on every pair of
inputs $(x, y)$.
\end{definition}

\begin{definition}[Protocol computes a function]
\lean{CommunicationComplexity.Deterministic.Protocol.Computes}
\statementsource{roughgarden-cc-bootcamp}{pdf-d54d402b16e2-p008-b002}
\statementsource{raoyehudayoff-ch01-deterministic-protocols}{pdf-fa052d8c95c0-p003-b003}
\label{CommunicationComplexity.Deterministic.Protocol.Computes}
\leanok
\uses{CommunicationComplexity.Deterministic.Protocol, CommunicationComplexity.Deterministic.Protocol.run}
A protocol $p$ \emph{computes} a two-argument function $f : X \to Y \to \alpha$ if
$p.\texttt{run}\,x\,y = f\,x\,y$ for all inputs $x : X$ and $y : Y$.
\end{definition}

\begin{definition}[Role swap]
\lean{CommunicationComplexity.Deterministic.Protocol.swap}
\label{CommunicationComplexity.Deterministic.Protocol.swap}
\leanok
\uses{CommunicationComplexity.Deterministic.Protocol}
Given a protocol $p$ over $X \times Y$, \texttt{swap} produces a protocol over $Y \times X$
by exchanging the roles of Alice and Bob throughout: every alice node becomes a bob node
and vice versa, while output nodes are left unchanged.
\end{definition}

\begin{theorem}[Role swap preserves protocol output]
\lean{CommunicationComplexity.Deterministic.Protocol.swap_run}
\label{CommunicationComplexity.Deterministic.Protocol.swap_run}
\leanok
\uses{CommunicationComplexity.Deterministic.Protocol, CommunicationComplexity.Deterministic.Protocol.run, CommunicationComplexity.Deterministic.Protocol.swap}
\difficulty{2}
Let $p$ be a deterministic two-party communication protocol with Alice's inputs drawn
from $X$, Bob's inputs drawn from $Y$, and outputs in $\alpha$, and let $x : X$ and $y :
Y$ be inputs. Then executing the role-swapped protocol on the input pair $(y, x)$
returns the same value as executing $p$ on $(x, y)$.
\end{theorem}

\begin{theorem}[Role swap preserves communication complexity]
\lean{CommunicationComplexity.Deterministic.Protocol.swap_complexity}
\label{CommunicationComplexity.Deterministic.Protocol.swap_complexity}
\leanok
\uses{CommunicationComplexity.Deterministic.Protocol, CommunicationComplexity.Deterministic.Protocol.complexity, CommunicationComplexity.Deterministic.Protocol.swap}
\difficulty{2}
Let $p$ be a deterministic two-party communication protocol with Alice's inputs drawn
from a type $X$, Bob's from a type $Y$, and outputs in a type $\alpha$. Let
$p^{\mathrm{swap}}$ be the protocol obtained by interchanging the roles of Alice and Bob
throughout — turning every Alice node into a Bob node and vice versa while leaving
output nodes unchanged. Then $p^{\mathrm{swap}}$ has the same communication complexity
as $p$: the worst-case number of bits exchanged is unaffected by the swap.
\end{theorem}

\begin{theorem}[Role swap is an involution]
\lean{CommunicationComplexity.Deterministic.Protocol.swap_swap}
\label{CommunicationComplexity.Deterministic.Protocol.swap_swap}
\leanok
\uses{CommunicationComplexity.Deterministic.Protocol, CommunicationComplexity.Deterministic.Protocol.swap}
\difficulty{2}
Let $p$ be a deterministic two-party communication protocol with Alice's inputs drawn
from a type $X$, Bob's from a type $Y$, and outputs in a type $\alpha$. Applying the
role swap twice — first exchanging the roles of Alice and Bob throughout $p$, then
exchanging them again — returns the original protocol $p$ unchanged.
\end{theorem}

\begin{theorem}[Swapping an Alice-rooted protocol to a Bob-rooted one]
\lean{CommunicationComplexity.Deterministic.Protocol.alice_to_bob}
\label{CommunicationComplexity.Deterministic.Protocol.alice_to_bob}
\leanok
\uses{CommunicationComplexity.Deterministic.Protocol, CommunicationComplexity.Deterministic.Protocol.complexity, CommunicationComplexity.Deterministic.Protocol.run, CommunicationComplexity.Deterministic.Protocol.swap, CommunicationComplexity.Deterministic.Protocol.swap_complexity, CommunicationComplexity.Deterministic.Protocol.swap_run}
\difficulty{1}
Let $X$, $Y$, and $\alpha$ be types, let $f : X \to \mathrm{Bool}$, and let $P$ assign
to each bit $b \in \mathrm{Bool}$ a deterministic communication protocol $P(b)$ over
input types $X$ (Alice) and $Y$ (Bob) with output type $\alpha$. Consider the protocol
whose root is an Alice node: on inputs $x$ and $y$ it applies $f$ to $x$ and continues
with $P(f(x))$. Then there exists a protocol $q$ over input types $Y$ (Alice) and $X$
(Bob) with output type $\alpha$ such that, for all $x \in X$ and $y \in Y$, executing
$q$ on inputs $y$ and $x$ returns the same value of $\alpha$ as executing the original
Alice-rooted protocol on inputs $x$ and $y$, and moreover $q$ has the same communication
complexity as that protocol.
\end{theorem}

\begin{theorem}[Turning a Bob-rooted protocol into a role-swapped protocol]
\lean{CommunicationComplexity.Deterministic.Protocol.bob_to_alice}
\label{CommunicationComplexity.Deterministic.Protocol.bob_to_alice}
\leanok
\uses{CommunicationComplexity.Deterministic.Protocol, CommunicationComplexity.Deterministic.Protocol.complexity, CommunicationComplexity.Deterministic.Protocol.run, CommunicationComplexity.Deterministic.Protocol.swap, CommunicationComplexity.Deterministic.Protocol.swap_complexity, CommunicationComplexity.Deterministic.Protocol.swap_run}
\difficulty{1}
Let $X$ and $Y$ be the input types of Alice and Bob and let $\alpha$ be the output type.
Fix a function $f : Y \to \mathrm{Bool}$ together with a family of deterministic
communication protocols $P(\mathtt{false})$ and $P(\mathtt{true})$, and let $R$ be the
protocol whose root is a bob node: on inputs $x$ (Alice) and $y$ (Bob) it computes the
bit $f(y)$ and then runs $P(f(y))$. Then there is a deterministic communication protocol
$q$ with the roles of Alice and Bob interchanged — Alice's input drawn from $Y$ and
Bob's from $X$ — such that for all $x$ and $y$, running $q$ on Alice-input $y$ and
Bob-input $x$ yields the same output value as running $R$ on Alice-input $x$ and
Bob-input $y$, and moreover $q$ has the same communication complexity as $R$.
\end{theorem}

\begin{definition}[Input pull-back]
\lean{CommunicationComplexity.Deterministic.Protocol.comap}
\label{CommunicationComplexity.Deterministic.Protocol.comap}
\leanok
\uses{CommunicationComplexity.Deterministic.Protocol}
Given functions $f_X : X' \to X$ and $f_Y : Y' \to Y$, \texttt{comap} pulls back a protocol
$p$ over $X \times Y$ to a protocol over $X' \times Y'$ by pre-composing every message
function with $f_X$ or $f_Y$ as appropriate, leaving output nodes unchanged.
\end{definition}

\begin{theorem}[Pull-back commutes with protocol execution]
\lean{CommunicationComplexity.Deterministic.Protocol.comap_run}
\label{CommunicationComplexity.Deterministic.Protocol.comap_run}
\leanok
\uses{CommunicationComplexity.Deterministic.Protocol, CommunicationComplexity.Deterministic.Protocol.comap, CommunicationComplexity.Deterministic.Protocol.run}
\difficulty{2}
Let $p$ be a deterministic two-party communication protocol with Alice's inputs drawn
from $X$ and Bob's from $Y$, producing values in $\alpha$, and let $f_X : X' \to X$ and
$f_Y : Y' \to Y$ be arbitrary maps of input types. Then for all $x' \in X'$ and $y' \in
Y'$, running the pull-back of $p$ along $f_X$ and $f_Y$ on the inputs $x'$ and $y'$
yields the same output as running $p$ itself on the transported inputs $f_X(x')$ and
$f_Y(y')$.
\end{theorem}

\begin{theorem}[Pull-back preserves communication complexity]
\lean{CommunicationComplexity.Deterministic.Protocol.comap_complexity}
\label{CommunicationComplexity.Deterministic.Protocol.comap_complexity}
\leanok
\uses{CommunicationComplexity.Deterministic.Protocol, CommunicationComplexity.Deterministic.Protocol.comap, CommunicationComplexity.Deterministic.Protocol.complexity}
\difficulty{2}
Let $p$ be a deterministic two-party communication protocol with Alice's inputs drawn
from $X$, Bob's inputs from $Y$, and outputs in $\alpha$, and let $f_X : X' \to X$ and
$f_Y : Y' \to Y$ be maps re-indexing the two input types. Then the pull-back protocol
over $X'$ and $Y'$ obtained by pre-composing each of Alice's message functions with
$f_X$ and each of Bob's message functions with $f_Y$ has the same communication
complexity as $p$.
\end{theorem}
